\section{The wallet and the PWA}\label{sec:wallet}

\subsection{Keys and the vault}

A wallet is a 12-word BIP-39 mnemonic. The full 64-byte BIP-39 seed is
derived through BIP-32 to
\begin{equation}
m/84'/\mathit{coin}'/0',
\end{equation}
with coin type 0 on mainnet and 1 elsewhere. The receive script is fixed at
$m/84'/\mathit{coin}'/0'/0/0$ and change addresses are
$m/84'/\mathit{coin}'/0'/1/i$; both are P2WPKH. The fixed receive script is a
deliberate v8 simplicity and recovery tradeoff. The watch set is fixed,
paginated to exhaustion, and a restored wallet reuses the same script;
address rotation is a future feature, not a consensus rule.

The ordinals branch is different. The branch is derived at
\begin{equation}
m/84'/\mathit{coin}'/0'/2/0.
\end{equation}
The branch point is normalized to even $Y$ and tweaked once with an empty
merkle root (the BIP86 \cite{bip86} tweak). The resulting output key is the
32-byte \code{owner}: it is the owner in the mint commit leaf and the commit
output's internal key, the owner committed by the four state leaves, the key
that signs crossings and transfers, and the wallet's current frog receive
address. The untweaked branch key is not itself an owner. The released
implementation tweaked the raw point without normalizing; a recovery scan
also watches that legacy address and derives the legacy state scripts.
Migration is explicit: a \code{migrateLegacyFrog} transfer moves the sat to
the current address. The old path is never rewritten in place. This key is
not used for funding. Coins at the frog address are never candidates for coin
selection, because spending one as fee would burn a frog.

The vault stores the BIP-39 entropy sealed with AES-256-GCM under a
non-extractable WebCrypto key-encryption key held in IndexedDB. Creation
writes, reads back and verifies the ciphertext before reporting success; a
storage denial is an error, never a ``sealed'' vault. The public ledger
(addresses, UTXOs, frog records, history) is stored separately. The ledger is
re-derivable from the chain; the 12 words are the only backup.

\begin{table}[htbp]
\centering
\small
\caption{Wallet derivation paths.}
\label{tab:derivation}
\begin{tabularx}{\linewidth}{@{}p{0.36\linewidth}p{0.16\linewidth}X@{}}
\toprule
\textbf{Path} & \textbf{Script} & \textbf{Purpose} \\
\midrule
$m/84'/\mathit{coin}'/0'/0/0$ & P2WPKH & Receive (fixed in v8) \\
$m/84'/\mathit{coin}'/0'/1/i$ & P2WPKH & Change \\
$m/84'/\mathit{coin}'/0'/2/0$ & P2TR & Frog receive and state owner \\
\bottomrule
\end{tabularx}
\end{table}

\subsection{Chain access}

The wallet has two HTTP ports, chosen by URL.

\paragraph{Esplora.} For an \code{https://} endpoint, the wallet uses the Esplora API \cite{esplora}:
\code{/blocks/tip/height}, \code{/address/{a}/utxo},
\code{/address/{a}/txs}, \code{/address/{a}/txs/chain[/{after_txid}]},
\code{/tx/{txid}/hex}, \code{/blocks/tip/hash},
\code{/block-height/{h}}, \code{/fee-estimates}, and \code{POST /tx}.
Address history is paginated to exhaustion; the first 50 rows are never
treated as complete.

\paragraph{Electrum over WebSocket.} For a \code{ws://} or \code{wss://}
endpoint, the wallet speaks JSON-RPC 2.0 and uses
\code{blockchain.scripthash.*}, \code{blockchain.transaction.get},
\code{blockchain.block.header(s)}, \code{blockchain.estimatefee}, and
\code{blockchain.transaction.broadcast}. A relay (\code{frogger-relay}) can
bridge raw TCP.

Discovery comes from the HOPPER indexer
(\code{/v1/froggers}, \code{/v1/frogger/{sat}}, \code{/v1/utxo/...}) or,
without an indexer, from stock ord (\code{/r/utxo/...},
\code{/content/{id}}, \code{/r/blockheight}). Ranges are mandatory before a
coin may be spent. An unattributable coin still counts in the balance, but it
is displayed as unspendable rather than spent blindly.

\paragraph{Tip catch-up.} Address scans are gated on the tip
$(\mathit{height},\mathit{hash})$: outside the catch-up window below, an
unchanged tip with no wallet transaction in flight costs no address queries.
The tip endpoint can publish a block before the address index has that
block's transactions. After an automatic pass that observes a new tip, the
wallet must therefore run two further full passes at
\code{TIP_RECHECK_MS} (about 1.2 s) even though the tip is unchanged. Those
passes bypass the gate and read the addresses again; only then may the gate
skip them. Without this grace window, the first pass could record the new tip
while holding the old UTXO set and the gate would keep the stale view until
the next block or a forced \code{RESCAN}. A manual \code{RESCAN} is a forced
pass; it does not itself create this automatic grace window.

\paragraph{Scan repaint.} A completed scan is also a UI repaint event. If it
lands while a foreground action holds the busy overlay, the repaint must be
queued and applied when that action releases the UI, not dropped. Otherwise
the Road or Pond can stay on the previous snapshot until the next block or a
manual \code{RESCAN} even though the wallet state is current.

A mint reveal is also self-describing: its tapscript witness contains the
canonical v8 body (or the v7 body for the grandfather), and output 0 pays the
wallet's ordinals address. A wallet
may recover an idle mint record from that witness when neither an indexer
nor ord answers, verifying the seed, sat and full body before trusting it.
This recovers mints paid directly to the wallet's ordinals address. It does
not recover a frog transferred in later, whose new output carries no mint
body; that discovery requires the indexer, or a working ord range view and a
record source.

\subsection{Screens and interaction}

The interface is a single-page PWA with the following screens:

\begin{center}
\small
\begin{tabular}{@{}llll@{}}
\toprule
Wallet & History & The Pond & The Road \\
Mint & Leaderboard & Detail & Bio \\
\bottomrule
\end{tabular}
\end{center}

\noindent
These are supplemented by onboarding (choose, show, confirm, recover,
mistake), the intro, an eight-page rules walkthrough, the $\equiv$ drawer,
send and receive sheets, the address book, backup words, a QR scanner, the
dead-frog page, and settings.

The visual language is fixed: dark lily-pad green \texttt{\#18241a}, card
\texttt{\#223226}, accent \texttt{\#f0c83c}, 8-bit monospace, square corners,
bark-framed buttons, lily-pad-bordered floating windows, and 16$\times$16
deterministic frog sprites. The art is generated from the seed, not fetched
from a server.

The road screen asks one question: \textbf{LEFT}, \textbf{STRAIGHT} or
\textbf{RIGHT}. The chosen button is the direction committed in the crossing
body and the leaf spent; the preview names it, and it cannot change after
approval. The screen states the rule in one line: the next block picks one
direction, and a Hopper that picked it dies. A frog's card shows its
generation badge (\code{G1}, \code{G2}, \ldots), its crossing count beside
the road mark, and \code{RIP} before the name once it is dead. Each settled
hop gets one road-log row. A crossed row says \textbf{survived}; it says
\textbf{came home} only when that settle's \code{immortal_height} is the
crossing height for the row ($\mathit{row.height}-\kappa$). A dead row says
\textbf{died}. If the frog record is absent, for example after a transfer
out of this wallet, the row is conservative and says \textbf{survived}
rather than claiming \code{HOME}.

\begin{figure}[htbp]
\centering
\includegraphics[width=0.82\linewidth]{figures/spritesheet.png}
\caption{Deterministic 16$\times$16 frog sprites. The same seed produces the
same sprite in the Rust implementation, the JavaScript wallet and the
inscribed fold.}
\label{fig:sprites}
\end{figure}

\subsection{Sound}

The score is one MIDI file (\code{music.mid}) played by \code{sound.js}. The
player parses the file once, maps ticks to seconds through the tempo map,
schedules notes with a 400 ms lookahead, and loops at the score's length.
Each note is one periodic-wave voice with a harmonic stack falling as
$1/k^{1.15}$, a 4 ms attack, and a pitch-dependent decay. The signal chain is
a 45 Hz highpass, a $+2.5$ dB shelf at 1.5 kHz, a compressor at $-16$ dB with
a 3.5:1 ratio, and a master gain. The gain is $0.5\cdot v^2$ for volume $v$.
Mute wins over volume, and nothing sounds before a user gesture.

\subsection{Emergency descriptor export}

The seed words are the only backup, but they do not describe how a custody
output is built. The wallet therefore offers \textbf{EXPORT EMERGENCY
DESCRIPTORS} in settings: one plain-text JSON document (also delivered as a
QR) that is enough for another taproot-capable tool to spend every staged
address without this wallet.

A \emph{staged address} is the P2TR state output of a frog. Per frog the
export carries:
\begin{center}
\begin{minipage}{0.92\linewidth}
\begin{lstlisting}[language={},caption={Per-frog emergency export fields.},label={lst:emergency}]
sat, gen, dist, state, name
mint_id            "<txid>i<index>"
state_outpoint     "<txid>:<vout>"      (when known)
state_value        sats
owner_xonly        32-byte hex
internal_key       50929b74...3ac0
output_key         32-byte x-only output key of the staged address
leaves[]           { dir, script_hex, leaf_hash, control_block }
send_leaf          { script_hex, leaf_hash, control_block }
descriptor         tr(<internal>,{raw(<L>),raw(<S>),raw(<R>),raw(<X>)})#checksum
\end{lstlisting}
\end{minipage}
\end{center}
A Hopper on a legacy path exports the released owner; a v7-wire Hopper also
uses the \code{FRG7CROSS} leaves. \code{owner_xonly} is that legacy owner.
Before displaying the document the wallet rebuilds each output key and
compares it with the scanned state scriptPubKey; a mismatch removes that frog
from the export.

and once for the wallet:
\begin{center}
\begin{minipage}{0.92\linewidth}
\begin{lstlisting}[language={},caption={Wallet-level emergency export fields.},label={lst:emergency-wallet}]
type               "hopper-v8-emergency"
version            1
network            mainnet | testnet | signet | regtest
fingerprint        BIP-32 master fingerprint (8 hex)
owner_path         m/84'/<coin>'/0'/2/0
descriptor         tr(<internal>,{raw(<L>),raw(<S>),raw(<R>),raw(<X>)})#checksum
                   (present only when exactly one staged tree; null otherwise)
generated_at       ISO-8601 UTC
tip_height         the height the export was built at
\end{lstlisting}
\end{minipage}
\end{center}

Each \code{descriptor} line uses the BIP-380 checksum and the miniscript
\code{raw()} fragment for that frog's four leaves. Many signers cannot parse
raw leaves, so the \code{leaves[]} hex and control blocks are the
authoritative form; the per-frog descriptor is a convenience for tools that
can. The wallet-level descriptor is present only when exactly one tree is
staged; with more than one it is \code{null}, and the per-frog descriptors
are the complete set.

Rules:
\begin{itemize}[leftmargin=1.4em,itemsep=0.25em]
  \item no seed words, no private keys, no xprv. The export is a \emph{layout}
        backup, not a key backup: whoever holds it still needs the seeds;
  \item the signing key is the BIP-32 key at \code{owner_path}, tweaked once
        with \code{TapTweak(P)} and an empty merkle root (the BIP86 tweak);
        the leaf's \code{owner_xonly} is the resulting output key. Spending
        uses leaf version \code{0xc0} and SIGHASH\_ALL (\code{0x01}), with the
        script and control block from \code{leaves[]};
  \item a receiving tool must verify the export before trusting it: rebuild
        each leaf hash, each branch, the root, the output key and every
        control block from the export's own bytes, and compare the result
        with the on-chain scriptPubKey of \code{state_outpoint}. The wallet
        runs the same check before it displays the export;
  \item the export is public data about the wallet's holdings. It is shown
        behind a confirmation and is never logged, uploaded, or put in a URL;
  \item the export is a snapshot, and every hop, prepare or transfer changes
        a state outpoint. The wallet stamps \code{tip_height} and regenerates
        the export after every action; a tool that sees a state outpoint
        already spent must discard the entry rather than guess.
\end{itemize}

\subsection{Wallet safety invariants}

These are mandatory before any transaction is signed or broadcast.

\begin{enumerate}[leftmargin=1.7em,itemsep=0.25em]
  \item \textbf{Complete ranges.} Every funding or state input has non-empty
        ordinal ranges whose lengths sum to its value. An empty
        ord/indexer answer is not proof that a coin has no inscription.
  \item \textbf{FIFO assignment.} The wallet recomputes the sat ranges of
        the exact signed outputs. A transfer, \textsc{prepare} or hop must
        place the Hopper sat in output 0 at offset 0; a send must not place
        any known protected sat in a recipient output, change, or the fee.
        RBF and the approval path repeat this proof.
  \item \textbf{Asset recheck.} Hop, prepare, transfer, send and mint
        revalidate every live input immediately before broadcast. The old
        exemption for send and mint is removed.
  \item \textbf{State discovery.} A state script alone never proves its sat
        ranges. Local discovery asks ord or the indexer for complete ranges
        and marks the coin unspendable until they match the value and
        contain the frog sat.
  \item \textbf{Recovery is fail-closed.} A restore paginates address
        histories to exhaustion, treats \code{more} or \code{null} as
        incomplete, reads at most a bounded checkpoint, and never
        synthesizes a range or a mint sat from a document.
  \item \textbf{Reorg identity.} Scan gates and outcome-hash caches are
        keyed by $(\mathit{height},\mathit{tip\ hash})$, not height alone.
        A same-height hash change forces a fresh pass.
  \item \textbf{Vault durability.} Vault creation writes, reads back and
        verifies the ciphertext; a storage failure is an error. Refresh
        summaries carry public keys only.
  \item \textbf{Service bounds.} Public reads are rate-limited, default
        300 per minute per caller; diagnostic \code{/v1/state/root} is not a
        hot path; and no state lock is held across an outbound ord call.
  \item \textbf{Scan freshness.} A new automatic tip gets two
        gate-bypassing catch-up passes (\code{TIP_RECHECK_MS}, about 1.2 s)
        so a lagging address index cannot freeze the UTXO view. A completed
        scan repaints the current view; if a foreground action owns the
        overlay, the repaint is queued and applied on release, never dropped.
\end{enumerate}

\subsection{The user's trust boundary}

The wallet never sends the mnemonic or the sealed entropy to a server. It
signs locally. It refuses a mint or a hop if the indexer's
\code{params_hash} differs from the embedded one, if the indexer does not
answer the parameter gate, or if its \code{network} does not match the
wallet's. It refuses a non-P2TR destination for custody. It refuses to spend
a coin whose ordinal ranges it cannot attribute. It never replaces a HOPPER
action with RBF. These refusals are part of the wallet's safety model, not UI
preferences.

\begin{invariant}{The wallet is not a custodian}
A wallet can be lost, replaced, or run offline. The frog remains a UTXO on
Bitcoin. Recovery is the 12 words; the public ledger is a cache.
\end{invariant}
